\documentclass[a4paper, leqno, 11pt]{article}
\usepackage[margin=1.2in]{geometry}
\usepackage[utf8]{inputenc}
\usepackage[T1]{fontenc}
\usepackage[english]{babel}
\usepackage{amsthm}
\usepackage{amssymb}
\usepackage{amsmath}
\usepackage{amsfonts}
\usepackage{mathtools}
\usepackage{latexsym}
\usepackage{graphicx}
\usepackage{float}
\usepackage{etoolbox}
\usepackage{hyperref}
\usepackage{subcaption}
\usepackage{tikz}
\usepackage{nicematrix}
\usepackage{booktabs}
\usepackage{pgfplots}
\pgfplotsset{compat=1.18} 
\usepackage{lipsum}
\usepackage{algorithm}
\usepackage{algpseudocode}
\usepackage{nccmath}
\setlength {\marginparwidth }{2cm}
\usepackage{todonotes}
\usepackage[most]{tcolorbox}
\usepackage{wrapfig}
\newtcolorbox[auto counter]{problem}[1][]{%
    enhanced,
    breakable,
    colback=white,
    colbacktitle=white,
    coltitle=black,
    fonttitle=\bfseries,
    boxrule=.6pt,
    titlerule=.2pt,
    toptitle=3pt,
    bottomtitle=3pt,
    title=GitHub repository of this project}

\RequirePackage[activate={true,nocompatibility},final,tracking=true,kerning=true,spacing=true]{microtype}
\SetTracking{encoding={*}, shape=sc}{40} % Spacing for Small Caps
\usepackage{dsfont}
\usepackage{setspace}
\usepackage{iftex}
% Define the new conditional
\newif\ifxetexorluatex

% Set it to true if compiling with XeTeX or LuaTeX
\ifxetex
  \xetexorluatextrue
\fi
\ifluatex
  \xetexorluatextrue
\fi

\ifxetexorluatex
  \usepackage{unicode-math}
  \setmainfont{EB Garamond}
  \setmathfont{Garamond Math}

  % Load some missing symbols from another font.
  \setmathfont{STIX Two Math}[%
    range = {
      \sharp,
      \natural,
      \flat,
      \clubsuit,
      \spadesuit,
      \checkmark
    }
  ]
  \setmonofont[Scale=MatchLowercase]{Source Code Pro}
\else
  \usepackage[lf]{ebgaramond}
  \usepackage[oldstyle,scale=0.7]{sourcecodepro}
  \singlespacing
\fi

\usepackage{tensors}


\renewcommand{\mathbb}[1]{\mathds{#1}}
\newcommand{\R}{\mathbb{R}}
\newcommand{\N}{\mathbb{N}}
\newcommand{\Z}{\mathbb{Z}}
\newcommand{\Q}{\mathbb{Q}}
\newcommand{\C}{\mathbb{C}}
\newcommand{\V}{\mathcal{V}}
\newcommand{\E}{\mathcal{E}}
\newcommand{\G}{\mathcal{G}}
\renewcommand{\L}{\mathbb{L}}

\newcommand{\defeq}{\vcentcolon=}
\newcommand{\eqdef}{=\vcentcolon}
\newcommand{\norm}[1]{\left\lVert#1\right\rVert}
\newcommand{\Kr}[3]{\mathcal{K}_{#1}(#2,#3)}

\makeatletter
\let\Vec\relax
\DeclareMathOperator{\Vec}{vec}
\let\vec\relax
\DeclareRobustCommand{\vec}{\Vec}
\DeclareMathOperator{\nnz}{nnz}
\makeatother

\newtheorem{theorem}{Theorem}[section]
\newtheorem{lemma}[theorem]{Lemma}
\newtheorem{proposition}[theorem]{Proposition}
\newtheorem{corollary}[theorem]{Corollary}
\theoremstyle{definition}
\newtheorem{definition}[theorem]{Definition}
\newtheorem{example}[theorem]{Example}
\theoremstyle{remark}
\newtheorem{remark}[theorem]{Remark}

% Define a distinct style for experiments
\newtheoremstyle{experimental}
  {\topsep}            % Space above
  {\topsep}            % Space below
  {\upshape}           % Body font: Upright (non-italic) suits empirical content
  {}                   % Indent amount
  {\bfseries} % Head font: Bold sans-serif for a modern, technical aesthetic
  {.}                  % Punctuation after theorem head
  {.5em}               % Space after theorem head
  {}                   % Theorem head spec

% Apply the style and define the unnumbered environment
\theoremstyle{experimental}
\newtheorem*{experiment}{Experiment}

\title{%
   Fast matrix multiplication using CP decomposition
    \\ \large Computational laboratory project report}
\author{Alberto Defendi\\Project supervisor: Leonardo Robol}

\date{}

\setlength{\parskip}{1em}
\setlength{\parindent}{0em}

\begin{document}
\maketitle

\begin{abstract}
    \noindent We develop a Rust
crate (library) that given matrices $ A,B \in \R^n$ computes their product $ C=AB $ using fast matrix multiplication algorithms, for any given CP decomposition.
\end{abstract}

{\setlength{\parskip}{0em}
\tableofcontents}

\section{Introduction}


\section{Notation and tensor preliminaries}

The relevant notation for our work is in Table~\ref{tab:notation}. We briefly review basic tensor preliminaries,
following the notation in \cite{BK25}.
A \emph{tensor} is a multidimensional array, and in this work we deal only with 3-dimensional (or 3-way)
tensors $ \Tn{x} \in \R^{m\times n\times p}$. The \emph{$ k $-th frontal slice} of $ \Tn{x} $ is the matrix
$ \Mx{X}_k $. For $ \Vc{u} \in \R^m, \Vc{v} \in \R^n, \Vc{w} \in \R^p $, we define the rank-1 \emph{outer
product} tensor $\Tn{x} = \Vc{u} \topp \Vc{v} \topp \Vc{w} \in \R^{m \times n \times p}$ with entries
$ x_{ijk} = u_iv_jw_k $. Addition of tensors is defined entry-wise. The \emph{rank} of a tensor $ \Tn{x} $
is defined as the minimum number of rank-one tensors that generate $ \Tn{x} $ as their sum.
We define the \emph{CP} decomposition of rank $ R $ of a tensor $\Tn{x}$ as the sum $ \Tn{x} =
\sum_{r=1}^R u_r \topp v_r \topp w_r $. We denote it by $ \Tn{x} = \dsqr{U,V,W} $, where $ U,V $
and $ W $ are matrices with $ R $ columns given by $ u_r,v_r $, and $ w_r $. An important operation
in the context of tensors is the \emph{mode$ -k $ product}, which is defined as. 
In this work, we will often use the equation $ \Tn{x} \ttm[1]
a^{\top} \ttm[2] b^{\top} = c $, with $ c_k = a^{\top}\Tm{x}{k}b = \sum_{i=1}^m\sum_{j=1}^n x_{ijk}$ to refer to matrix multiplication algorithms.

\begin{table}[tb]
\centering
\caption{Summary of notation.}
\label{tab:notation}
\begin{tabular}{ll}
\toprule
$\angb{M,K,N}$ & Bilinear form representing base case matrix multiplication \\
& of an $M \times K$ matrix by a $K \times N$ matrix. \\
$\msiz{P,Q,R}$ & Dimensions of actual matrices that are multiplied \\
& ($P \times Q$ matrix multiplied by $Q \times R$ matrix). \\
$\dsqr{\Mx{U}, \Mx{V}, \Mx{W}}$ & Factor matrices corresponding to a tensor CP \\
& decomposition that provides a fast algorithm. \\
$\Tn{X}$ & Order-3 tensor. \\
$\Tn{X} = \Vc{u} \topp \Vc{v} \topp \Vc{w}$ & Rank-1 tensor with entries $X_{ijk} = u_i v_j w_k$. \\
$\Tm{X}{i}$ & $i$th frontal slice of $\Tn{X}$, the matrix of entries $X_{:,:,i}$. \\
$\Vc{a}_k$, $A_{ij}$ & $k$th column and $i, j$ entry of matrix $\Mx{A}$. \\
$\vec(\Mx{A})$ & Row-order vectorization of the entries of $\Mx{A}$. \\
$\nnz(\cdot)$ & Number of non-zero entries of an object. \\
\bottomrule
\end{tabular}
\end{table}


\section{Fast matrix multiplication}

Matrix multiplication is a bilinear operation. \todo{explain}
We indicate with $ \angb{m,n,p} $ the bilinear form that represents the multiplication between a
matrix $ \Tm{A} \in \R^{m\times n} $ with $ \Tm{B} \in \R^{n\times p} $. Fixing the natural
ordering on the entries of $ \Tm{A} $ and $ \Tm{B} $ and the inverse ordering on $ \Tm{C} =
\Tm{A}\Tm{B} $, we get a representation of $ \Tn{x} \in \R^{mn\times np\times mp} $ of the matrix
multiplication $ \angb{m,n,p} $ such that
\begin{equation}
  \label{eq:mulvec}
\Vec(\Tm{C}^{\top}) = \Tn{X} \ttm[1] \Vec(\Tm{A})^{\top} \ttm[2] \Vec(\Tm{B})^{\top}
.\end{equation}
Note that we fix the linear order on $ A $ and $ B $, while we fix the reverse ordering on $ C $ by
applying the transpose operation.

\todo{Tensor definition}

If we are given a CP decomposition of rank $ r $ of the tensor $ \Tn{x}  = \dsqr{U,V,W}$, we can
rewrite Equation \ref{eq:mulvec} into
\begin{equation}
  \label{eq:mulvec:CP}
  \Vec(\Tm{C}^{\top}) = \sum_{l=1}^r (\underbrace{\Vc{u_l^{\top}} \Vec(\Tm{A})}_{s_l}) \cdot
  (\underbrace{\Vc{v_l^{\top}}\vec(\Tm{B})}_{t_l})\Vc{w_l} = \sum_{l=1}^r (\underbrace{s_l \cdot
  t_l}_{m_l})w_l,
\end{equation}
where $ \Vc{u_l}, \Vc{v_l}, \Vc{w_l}$ denote the columns of $ \Tm{U},\Tm{V} $ and $ \Tm{W} $. This
reduces the number of active multiplications to the rank $ r $ of $ \Tn{x} $.


\subsection{The classic matrix multiplication algorithm tensor formulation}

We present the $ \angb{2,2,2} $ classical matrix multiplication algorithm in the context of tensors.
When $ m,n,p >2 $ we can always partition them in blocks as:

\begin{equation}
  \label{eq:blocks}
  \begin{bNiceArray}{cc}[first-row, first-col, margin, xdots/horizontal-labels]
    & \Hbrace{1}{\lfloor n/2 \rfloor} & \Hbrace{1}{\lceil n/2 \rceil} \\
    \Vbrace{1}{\lfloor m/2 \rfloor} & A_{11} & A_{12} \\[2ex]
    \Vbrace{1}{\lceil m/2 \rceil}   & A_{21} & A_{22}
  \end{bNiceArray}, \qquad \qquad
  \begin{bNiceArray}{cc}[first-row, first-col, margin, xdots/horizontal-labels]
    & \Hbrace{1}{\lfloor p/2 \rfloor} & \Hbrace{1}{\lceil p/2 \rceil} \\
    \Vbrace{1}{\lfloor n/2 \rfloor} & B_{11} & B_{12} \\[2ex]
    \Vbrace{1}{\lceil n/2 \rceil}   & B_{21} & B_{22}
  \end{bNiceArray},
\end{equation}

where the sizes of the subblocks are depicted in the above equation. This algorithmic formulation is
often called \emph{divide et impera}, as it consists in breaking the problem into larger
subproblems, and later recombining the results.
In particular, we can view the product $ AB $ as the product between $ 2\times 2 $ block matrices 
\begin{equation}
  \label{eq:matprod}
  \begin{bmatrix}
    C_{11} & C_{12} \\
    C_{21} & C_{22}
  \end{bmatrix} =
  \begin{bmatrix}
    A_{11} & A_{12} \\
    A_{21} & A_{22}
  \end{bmatrix} \cdot 
  \begin{bmatrix}
    B_{11} & B_{12} \\
    B_{21} & B_{22}
  \end{bmatrix} =
  \begin{bmatrix}
    A_{11}B_{11} + A_{12}B_{21} & A_{11}B_{12} + A_{12}B_{22} \\
    A_{21}B_{11} + A_{22}B_{21} & A_{21}B_{12} + A_{22}B_{22}
  \end{bmatrix}.
\end{equation}
The naive algorithm proceeds by combining a set of eight matrix multiplications with four
additions:
% Intermediate Matrices
% Intermediate Matrices
\[
\begin{aligned}
M_1 &= A_{11} \cdot B_{11} &\quad M_2 &= A_{12} \cdot B_{21} &\quad M_3 &= A_{11} \cdot B_{12} \\
M_4 &= A_{12} \cdot B_{22} &\quad M_5 &= A_{21} \cdot B_{11} &\quad M_6 &= A_{22} \cdot B_{21} \\
M_7 &= A_{21} \cdot B_{12} &\quad M_8 &= A_{22} \cdot B_{22} & &
\end{aligned}
\]

% Resulting Matrices
\[
\begin{aligned}
C_{11} &= M_1 + M_2 &\qquad C_{12} &= M_3 + M_4 \\
C_{21} &= M_5 + M_6 &\qquad C_{22} &= M_7 + M_8
\end{aligned}
\]

The number of flops performed by standard matrix multiplication for $ N\times N $ matrices is $ T(N)
= 8T(\frac{N}{2}) + 4(\frac{N}{2})^2$ for $ N>1 $. This can be solved by the master method to get
$ T(N)  = 2N^3 - N^2 $, where we assume that N is a power of two. In Section \ref{sec:dynpeel} we
explain how to handle all dimensions.

For example, when $ m=n=p=2 $, we get the tensor with frontal slices
\[
\Mx{X}_1 = \begin{bmatrix}
  1 & 0 & 0 & 0 \\
  0 & 0 & 0 & 0 \\
  0 & 1 & 0 & 0 \\
  0 & 0 & 0 & 0
\end{bmatrix}, \quad
\Mx{X}_2 = \begin{bmatrix}
  0 & 0 & 1 & 0 \\
  0 & 0 & 0 & 0 \\
  0 & 0 & 0 & 1 \\
  0 & 0 & 0 & 0
\end{bmatrix}, \quad
\Mx{X}_3 = \begin{bmatrix}
  0 & 0 & 0 & 0 \\
  1 & 0 & 0 & 0 \\
  0 & 0 & 0 & 0 \\
  0 & 1 & 0 & 0
\end{bmatrix}, \quad
\Mx{X}_4 = \begin{bmatrix}
  0 & 0 & 0 & 0 \\
  0 & 0 & 1 & 0 \\
  0 & 0 & 0 & 0 \\
  0 & 0 & 0 & 1
\end{bmatrix}.
\]
This yields, for example,
\begin{equation*}
    X_3 \ttm[1] \Vec(A)^{\top} \ttm[2] \Vec(B)^{\top} =
    \begin{bmatrix}
      a_{11} \; a_{21} \; a_{12} \; a_{22}
    \end{bmatrix}
    \begin{bmatrix}
      0 & 0 & 0 & 0 \\
      1 & 0 & 0 & 0 \\
      0 & 0 & 0 & 0 \\
      0 & 1 & 0 & 0
    \end{bmatrix}
    \begin{bmatrix} b_{11} \\ b_{21}\\b_{12}\\b_{22} \end{bmatrix} 
    = a_{21}b_{11} + a_{22}b_{21} = c_{21}.
\end{equation*}

Note that the formula in Equation~\ref{eq:mulvec:CP} in the context of block matrix multiplication in
Equation~\ref{eq:matprod} yelds to
\begin{equation*}
  \Vec(A) = \begin{bmatrix} A_{11}\\A_{21}\\A_{12}\\A_{22} \end{bmatrix} 
, \quad
\Vec(B) = \begin{bmatrix} B_{11}\\B_{21}\\B_{12}\\B_{22} \end{bmatrix},\quad
\Vec(C^{\top}) = \begin{bmatrix} C_{11}\\C_{12}\\C_{21}\\C_{22} \end{bmatrix}.
\end{equation*}
Also the factors $ s_l $ and $ t_l $ in Equation~\ref{eq:mulvec:CP} become
\begin{equation*}
  s_l = u_l^{\top}\Vec(A) = \sum_{i=1}^4 U_{li}\Vec(A)_i, \quad t_l = V_l^{\top}\Vec(B) =
  \sum_{i=1}^4 V_{li}\Vec(B)_i
\end{equation*}


\subsection{Strassen algorithm using CP}
\label{sec:cpstrassen}

The idea of fast matrix multiplication algorithms is to perform fewer recursive matrix
multiplications at the expense of more matrix additions. The most well known fast algorithm is due to Strassen, and follows the same block structure:

\begin{align*}
S_1 &= A_{11} + A_{22} & S_2 &= A_{21} + A_{22} & S_3 &= A_{11} \\
S_4 &= A_{22} & S_5 &= A_{11} + A_{12} & S_6 &= A_{21} - A_{11} \\
S_7 &= A_{12} - A_{22} \\
T_1 &= B_{11} + B_{22} & T_2 &= B_{11} & T_3 &= B_{12} - B_{22} \\
T_4 &= B_{21} - B_{11} & T_5 &= B_{22} &T_6 &= B_{11} + B_{12} \\
T_7 &= B_{21} + B_{22}
\end{align*}
\begin{align*}
M_r &= S_rT_r, \quad 1 \le r \le 7 \\
C_{11} &= M_1 + M_4 - M_5 + M_7 & C_{12} &= M_3 + M_5 \\
C_{21} &= M_2 + M_4 & C_{22} &= M_1 - M_2 + M_3 + M_6
\end{align*}

This algorithm uses $ 7 $ matrix multiplications and $ 18 $ matrix additions. The number of flops
performed by the algorithm is $ T(N) = 7T(\frac{N}{2}) + 18(\frac{N}{2})^2$ for $ N>1 $, solving the
recursion the complexity is $ T(N) = 7N^{\log_2 7} - 6N^2 = O(N^{\log_2 7}) = O(N^{2.81}) $.

\todo{complete}


There are natural extensions to Strassen algorithm:

\begin{enumerate}
  \item We can't use less than 7 number of additions for $ \angb{2,2,2} $.
  \item Reduce the number of addition from 18 down to 15, which leads to Strassen-Winograd algorithm
    \todo{cite}.
  \item We can change constant on the leading term by chosing a bigger base case dimension (and
    using the standard algorithm (also called \emph{gemm}) on the base case).
  \item We can use blocking schemes different from $ \angb{2,2,2} $.
\end{enumerate}

\subsection{Fast algorithms as low-rank tensor decomposition}

Let $ X,Y,Z $ be finite-dimensional vector spaces of dimensions $ I,J,K $ respectively. A bilinear form is a mapping $ f : X \times Y \to
Z$ that is linear in each entry, in other words such that $ f(x, \cdot) $ and $ f(\cdot,y) $ are linear for all
$ x\in X $ and $ y\in Y $. Any bilinear form can be represented by a matrix $ M $ of coefficients: $
f(x,y) = x^{\top}My = \sum_i \sum_j m_{ij}x_iy_j $, and we can write it component-wise using a
tensor $ \Tn{t} $ of coefficients $ t_{ijk} $
\begin{equation}
  f_k(x,y) = \sum_{i=1}^I \sum_{j=1}^J t_{ijk}x_iy_j, \quad \text{ for all } 1\le k\le K,
\end{equation}
or in more succint tensor notation, $ f = \Tn{x} \ttm[1] x \ttm[2] y $.

\todo{Extend and move before defining matmul formula}

\section{Implementation and practical consideration}

\subsection{Algorithm schema}

We now discuss the algorithm schema for Strassen algorithm for clarity, though results generalize
easily for any $ \angb{m,n,p} $ base case. Algorithm \ref{alg:strassen} shows our implementation of
Strassen's algorithm, which easily extends to other CP algorithms by changing base conditions and
block sizes---this is how it is implemented in practice.

\begin{algorithm}[ht!]
\caption{Strassen algorithm using CP factorization.}
\label{alg:strassen}
\begin{algorithmic}[1]
\Procedure{Strassen}{$ A,B, N $}
\If{\Call{stop-recursing}{} is True}
    \State \Return{$ A\cdot B $}
  \EndIf
  \If{$ N = 2 $}
  \State \Call{Strassen-Base-Matmul}{$ A,B $}
  \EndIf
  \If{$ N \mod 2 > 0 $}
    \State \Call{dynamic-peeling}{$A,B,N$} \Comment{Recursive call on peeled blocks.}
  \EndIf
  \State Compute $ \vec(A) $ and $ \vec(B) $ with block size $ \frac{N}{2} $.
  \For{$ r = 1, \ldots, R $}
  \State  $ S_r \leftarrow \sum_{i=1}^{4} U_{i,r} \vec(A)_i$
  \State  $ T_r \leftarrow \sum_{i=1}^{4} V_{i,r} \vec(B)_i$
  \State $ M_r \gets\Call{Strassen}{S_r,T_r,\frac{N}{2}} $  \Comment{Recursive call.}
  \EndFor
  \For{$ r = 1,\ldots,R $}
    \For{$i = 1,\ldots,2$}
      \For{$j = 1,\ldots,2$}
        \State $ w \gets W_{\L^*(i,j)},r = W_{2i + j,r}$
        \If{$w \neq 0$}
          \State $ C_{i\cdot \frac{N}{2}:(i+1)\cdot \frac{N}{2}, j\cdot
          \frac{N}{2}:(j+1)\frac{N}{2}} \gets w \cdot M_r $
        \EndIf
      \EndFor
    \EndFor
  \EndFor
  \State \Return $ C $.
\EndProcedure
\end{algorithmic}
\end{algorithm}


\subsection{Toeplitz cluster}

The project experiments were executed on Toeplitz cluster on \todo{node specs}.

\subsection{Technical overview}

The implementation consists in a Rust porting of the c++ project made by Benson and Ballard summarized in
\cite{bensonballard2015}, available in the relative repository under BSD 2-Clause. They
provide a fast and efficient framework for benchmarking parallel and sequential fast matrix
multiplication algorithms, achieved trough code generation, using lapack and MKL dgemm for numerical
computation, and OpenMP for parallelization.
Our project differs from their implementation mainly for the language (Rust) and code
generation---we apply CP matrix multiplication without code generation, which delivers simpler code
at the cost of speed \todo{is this True?}. Also, our project uses the numerical library Faer (cite)
with Rust's Rayon parallelization,
and selectively switches to Faer or MKL Dgemm for matrix multiplication,
while the other uses Lapack and MKL Dgemm, and OpenMP for parallelization.


Following the work \cite{bensonballard2015} we implement a fast multiplication algorithm based on CP
decomposition given the $ U,V $and $ W $ matrices representing the algorithm.


Fast matrix multiplication algorithm such as Strassen face two significant problems in practical
implementation, which are recursion cutoff and input matrices sizes not matching with base case. The
former can solved by \ldots and the latter by dynamic peeling~\ref{sec:dynpeel}.

\todo{Here comparison with Ballard et Al (used for comparison)}


\subsection{Recursive cutoff strategies}

We implement two strategies for stopping recursion before calling the vendor library, which are by
recursion depth and by matrices size. The former stops when a fixed number $ L $ of recursion levels
is reached, and the latter when matrices sizes reaches a given size.

\todo{Cutoff comparisons in sequential case??}

\subsection{Dynamic Peeling}
\label{sec:dynpeel}

In Strassen algorithm we want to multiply $ \angb{m,n,p} $ matrices by splitting in four blocks as
in Equation~\ref{eq:blocks}. If $ m,n $ and $ p $ are even, then it is not possible to split blocks
in even sizes. The easiest fix for this is padding the matrix with zeros until an even size is
reached, but this adds a significant memory waste. A more efficient approach, called \emph{dynamic
peeling}, consists in stripping the extra row off and/or column and adding their contribution to the
final results. In detail,

For the matrix multiplication $C = AB$, where $A$ is an $m \times n$ matrix and $B$ is an $n \times
p$ matrix. The matrix $A$ is split into an $(m-1) \times (n-1)$ core matrix $A_{11}$, alongside its
peeled boundary vectors $a_{12}$, $a_{21}$, and the scalar corner element $a_{22}$. Similarly, the matrix $B$ is split into an $(n-1) \times (p-1)$ core matrix $B_{11}$, with its corresponding peeled components $b_{12}$, $b_{21}$, and $b_{22}$:


\begin{equation*}
  A = \begin{bNiceArray}{c|c}[first-row, first-col, last-col, margin, xdots/horizontal-labels]
    & \Hbrace{1}{n-1} & \Hbrace{1}{1} \\
    & & &\Vbrace{3}{m-1}  \\
    & A_{11} & a_{12} & \\
    & & & \\ \hline
    & a_{21} & a_{22} & \Vbrace{1}{1}
  \end{bNiceArray} \qquad,\qquad
  B = \begin{bNiceArray}{c|c}[first-row, first-col, last-col, margin, xdots/horizontal-labels]
    & \Hbrace{1}{p-1} & \Hbrace{1}{1} \\
    & & &\Vbrace{3}{n-1}  \\
    & B_{11} & b_{12} & \\
    & & & \\ \hline
    & b_{21} & b_{22} & \Vbrace{1}{1}
  \end{bNiceArray}
\end{equation*}

The final block matrix product is computed directly through a combination of the core Strassen multiplication ($A_{11}B_{11}$) and low-rank vector/scalar fixup modifications:
\begin{equation*}
\begin{bNiceArray}{c|c}[first-row, first-col, last-col, margin, xdots/horizontal-labels]
    & \Hbrace{1}{p-1} &\Hbrace{1}{1} \\
    & & &  \Vbrace{3}{m-1} \\
    & C_{11} & c_{12} & \\
    & & & \\ \hline
    & c_{21} & c_{22} & \Vbrace{1}{1} 
  \end{bNiceArray} \qquad =
\begin{bNiceArray}{c|c}[]
    &   \\
    A_{11}B_{11} + a_{12}b_{21} & A_{11}b_{12} + a_{12}b_{22} \\
    &   \\ \hline
     a_{21}B_{11} + a_{22}b_{21} & a_{21}b_{12} + a_{22}b_{22}
  \end{bNiceArray}
\end{equation*}
Here, only the sub-product $A_{11}B_{11}$ is executed recursively using Strassen's matrix
multiplication, while all other terms are computed via standard matrix multiplication.



\section{Parallel algorithms for shared memory}

Matrix multiplication is a problem subject to both computational and memory bound. The first is
limited by the amount of computation the device can handle (i.e cores,cpu/gpu speed), while the
latter by physical storage of the matrices (i.e available RAM memory). In a parallel implemntation,
when dealing with recursive algorithms which use divide-et-impera (like Strassen), it is important
in practice to balance algorithmic load between threads to avoid overhead caused by both taking
recursive branches and doing parallel matrix multiplication. To address this issue, we adopt the
scheduling schema suggested by Benson and Ballard in their work~\cite{bensonballard2015}.

\section{Finding the optimal recursive cutoff}

In divide et impera algorithms such as Strassen, we can't afford taking many recursive stes, as
spawning many children constitutes a significant recursive overhead.
Thus, finding an optimal cutoff is important for stopping the recursion and switch to standard gemm.
We try two techniques: the first is breaking by size, that is when the sizes of matrices in the recursion fall
under a certain size, the latter is taking a fixed number $ L $ of recursions.
The authors in~\cite{bensonballard2015} explain an euristical process for finding the optimal cutoff
size, which can be resumed in looking at the performance plot of the gemm function, where (figure)
the function exibits a ramp-up and then flatten for large dimensions. They take a
recursion step only if the recursive sub-problem fall in the flat part of the curve. This is done
manually by the authors. During the experimentation we developed a function to plot this curve, and
its derivative (using splines) to automatically find this size. In practice, we never use this as we
believe that our code is not efficient enough to implement this strategy.

\subsection{Implementation}

We use Rust's Rayon to handle parallelization, and we can summarize the parallel scheduling as following: 

\begin{itemize}
  \item \textbf{BFS:} Each recursive matrix multiplication routine and the associated matrix
    additions are launched by Rayon as an iterable, and then are collected in an array of matrices $
    \left[M_1, \ldots, M_r \right]$.
  \item \textbf{DFS:} Each vendor matrix multiplication call uses all thread, and matrix
    multiplications are always fully parallelized.
  \item \textbf{Hybrid:} The Hybrid approach developed in \cite{bensonballard2015} compensates the
    for the load imbalance in BFS by applying DFS on a subset of the base case problems. With $ L $
    levels of recursion and $ P $ threads, the hybrid algorithm applies task parallelism (BFS) to
    the first $ R^L - (R^L \mod P) $ multiplication. Because the number of subproblems is a multiple
    of $ P $, this part is load balanced. On the remaining $ R^L \mod P $ subproblems, all threads
    are used on each matrix multiplication (DFS).
\end{itemize}


\begin{figure}[ht]
  \centering
  \resizebox{0.9\textwidth}{!}{%
    \begin{tikzpicture}[
    font=\sffamily,
    % Define a style for the tree nodes
    rootNode/.style={text=red!80!black, font=\large\bfseries},
    leafNode/.style={text=black, font=\small},
    activeLeaf/.style={text=red!80!black, font=\small},
    % Step labels
    labelStyle/.style={font=\Large\bfseries\sffamily}
]

    % ==========================================================
    % BREADTH-FIRST-SEARCH (BFS) TREE - LEFT SIDE
    % ==========================================================
    \node[labelStyle] (bfsTitle) at (0, 2) {Breadth-First-Search (BFS)};

    % Root node
    \node[rootNode] (bfsRoot) at (0, 0) {$A \cdot B$};

    % Leaf nodes (All 7 multiplications executed in parallel)
    \node[activeLeaf] (b0) at (-3.3, -2.5) {$T_0 \cdot S_0$};
    \node[activeLeaf] (b1) at (-2.2, -2.5) {$T_1 \cdot S_1$};
    \node[activeLeaf] (b2) at (-1.1, -2.5) {$T_2 \cdot S_2$};
    \node[activeLeaf] (b3) at (0.0, -2.5) {$T_3 \cdot S_3$};
    \node[activeLeaf] (b4) at (1.1, -2.5) {$T_4 \cdot S_4$};
    \node[activeLeaf] (b5) at (2.2, -2.5) {$T_5 \cdot S_5$};
    \node[activeLeaf] (b6) at (3.3, -2.5) {$T_6 \cdot S_6$};

    % Connect lines (All branches highlighted in red for parallel processing)
    \foreach \leaf in {b0, b1, b2, b3, b4, b5, b6} {
        \draw[red!80!black, line width=1.5pt] (bfsRoot.south) -- (\leaf.north);
    }


    % ==========================================================
    % DEPTH-FIRST-SEARCH (DFS) TREE - RIGHT SIDE
    % ==========================================================
    \begin{scope}[xshift=9cm]
        \node[labelStyle] (dfsTitle) at (0, 2) {Depth-First-Search (DFS)};

        % Root node
        \node[rootNode] (dfsRoot) at (0, 0) {$A \cdot B$};

        % Leaf nodes (Multiplications executed sequentially, highlighting one step)
        \node[leafNode] (d0) at (-3.3, -2.5) {$T_0 \cdot S_0$};
        \node[leafNode] (d1) at (-2.2, -2.5) {$T_1 \cdot S_1$};
        \node[activeLeaf] (d2) at (-1.1, -2.5) {$T_2 \cdot S_2$}; % Active branch step
        \node[leafNode] (d3) at (0.0, -2.5) {$T_3 \cdot S_3$};
        \node[leafNode] (d4) at (1.1, -2.5) {$T_4 \cdot S_4$};
        \node[leafNode] (d5) at (2.2, -2.5) {$T_5 \cdot S_5$};
        \node[leafNode] (d6) at (3.3, -2.5) {$T_6 \cdot S_6$};

        % Connect lines (Inactive branches are black, active branch step is red)
        \foreach \leaf in {d0, d1, d3, d4, d5, d6} {
            \draw[black, line width=1.5pt] (dfsRoot.south) -- (\leaf.north);
        }
        \draw[red!80!black, line width=2pt] (dfsRoot.south) -- (d2.north);
    \end{scope}

\end{tikzpicture}

  }%
  \caption{BFS and DFS load balancing. The parallelized tasks are in red.}
  \label{fig:unfoldings}
\end{figure}


\section{Experimental results}

\subsection{Metrics}

All benchmarks are recorded using Rust's system clock, performing warm-up before doing measuring
each matrix multiplication. In order to compare the performance of matrix multiplication algorithms with different
ocmputational costs, we use the \emph{effective} GFLOPS metric for $ \angb{m,n,r} $ matrix
multiplication:
\begin{equation}
  \text{effective GFLOPS} = \frac{2mnp - mp}{\text{time in seconds} \cdot 10^9},
\end{equation}
where the numerator can be rewritten as $2mnp - mp = mnp + m(n-1)p$, which is the number of
multiplications and additions performed for classical matrix multiplication. This allows us to
compare all of the algorithms on an inverse-time scale, normalized by problem size
\cite{bensonballard2015}. For fast algorithms, considering both multiplications and additions in the metric is more
accurate than taking multiplications only \cite{commavoidingstrass}.


\subsection{Measuring impact of base gemm}

We aim to understand how the choice of the base matrix multiplication impacts performance of fast
algorithms. To so so, we compare Faer

\subsection{Comparison of MKL dgemm FFI and Faer matmul}

Though we choose mkl-dgemm for a better comparison, our ffi bindings easily extend to any other gemm, such as OpenBlas, which we leave to further work.

\subsection{Numerical comparison with Benson\&Ballard Strassen}


\clearpage
\bibliographystyle{unsrt}
\bibliography{ref}
\nocite{*}



\end{document}
